% ===========================================================================
%  dodatekQ_coarse_graining_formal.tex
%  Teoria Generowanej Przestrzeni (TGP) — wersja 1
%
%  DODATEK Q: Formalne sformułowanie gruboziarnienia substrat→Φ (OP-2)
%  -------------------------------------------------------------------
%  Status: PROGRAM BADAŃ — precyzyjne sformułowanie problemu OP-2
%
%  Cel niniejszego dodatku NIE jest zamknięcie OP-2,
%  lecz sformułowanie go z matematyczną precyzją, podanie warunków
%  wystarczających dla zamknięcia oraz wskazanie drogi numerycznej.
%
%  Stan wiedzy po sek10 i dodatkuN:
%   - K(φ)=K_geo·φ² wyprowadzone z substratu (sek10)
%   - Punkt stały WF istnieje w LPA (dodatekN)
%   - K(0)=0 chronione przez Z₂ (sek10 + dodatekN)
%   - BRAKUJE: formalna zbieżność Φ_B → Φ_ciągłe
% ===========================================================================

\section*{Dodatek Q: Formalne gruboziarnienie substrat $\Gamma \to \Phi$}
\addcontentsline{toc}{section}{Dodatek Q: Gruboziarnienie substrat→Φ}
\label{app:coarse_graining}

\subsection*{Q.1 Sformułowanie problemu}

Teoria TGP opiera się na trzech poziomach:
\begin{enumerate}
  \item \textbf{Substrat} $\Gamma = (V,E)$ z hamiltonianem
        $H_\Gamma = -J\sum_{\langle ij\rangle}(\phi_i\phi_j)^2$,
        zmiennymi $\phi_i \in \mathbb{R}$, symetrią $\mathbb{Z}_2$.
  \item \textbf{Pole ciągłe} $\Phi(x)$ spełniające równanie pola TGP
        (sek08, eq.~\eqref{eq:action}).
  \item \textbf{Geometria} $g_{\mu\nu}$ emergentna z $\Phi$.
\end{enumerate}
Przejście $1 \to 2$ jest formalnie nieokreślone. Mówimy, że:
\begin{equation}
  \Phi(x) = \frac{1}{N_B}\sum_{i \in B(x)} \phi_i^2,
  \qquad N_B = \bigl|\{i : x_i \in B(x)\}\bigr|,
  \label{eq:block_avg}
\end{equation}
gdzie $B(x)$ jest blokiem o liniowym rozmiarze $b$ zawierającym węzeł $x$.
\textbf{Problem OP-2}: wykazać, że $\Phi_B(x) \to \Phi(x)$ w granicy
$b/a_\mathrm{sub} \to \infty$ oraz że $\Phi(x)$ spełnia równanie pola TGP
w sensie słabym.

\subsection*{Q.2 Warunki wystarczające (program badań)}

\begin{definicja}[Operator blokowania]
Definiujemy operator blokowania $\mathcal{T}_b$ na przestrzeni funkcjonałów:
\begin{equation}
  \mathcal{T}_b[S][\Phi_B] = -\ln \int \prod_{i \in V}
    \frac{d\phi_i}{\sqrt{2\pi}}
    \prod_{B} \delta\!\bigl(\Phi_B - \tfrac{1}{N_B}{\textstyle\sum}_{i\in B}\phi_i^2\bigr)
    \, e^{-S[\{\phi_i\}]},
  \label{eq:blocking_operator}
\end{equation}
gdzie całkowanie biegnie po konfiguracji dyskretnej przy ustalonym $\Phi_B$.
\end{definicja}

\begin{hipoteza}[Atraktor operatora $\mathcal{T}_b$]
\label{hyp:attractor}
W przestrzeni funkcjonałów z normą $\|\cdot\|_H$ operator $\mathcal{T}_b$
jest odwzorowaniem zwężającym (contractor) w sensie Banacha dla $b > b_*$,
gdzie $b_*$ jest związane ze skalą korelacji $\xi_\mathrm{corr}$.
Jego unikalny punkt stały $S^*[\Phi]$ jest tożsamy z działaniem TGP
(sek08, eq.~\eqref{eq:action}).
\end{hipoteza}

\begin{uwaga}[Status Hipotezy~\ref{hyp:attractor}]
Hipoteza~\ref{hyp:attractor} jest \emph{mocną} wersją OP-2.
Słabsza wersja (wystarczająca dla fizyki TGP) to jedynie zbieżność
$\mathcal{T}_b \to S^*$ w sensie dystrybucyjnym bez dowodu
zwężania. Dowód pełny jest otwartym problemem matematycznym.
\end{uwaga}

\subsection*{Q.3 Trzy argumenty za zbieżnością}

\paragraph{Argument Q-I: Teoria ERG (numeryczny).}
Z Dodatku~\ref{app:ERG} (LPA Wettericha):
\begin{itemize}
  \item Punkt stały Wilsona–Fishera istnieje i jest \emph{unikalny}
        w klasie teorii z symetrią $\mathbb{Z}_2$, $d=3$, $n=1$.
  \item Warunek $K(0) = 0$ jest chroniony przez $\mathbb{Z}_2$
        podczas całego przepływu.
  \item $K_\mathrm{IR}/K_\mathrm{UV} \approx 1{,}13$:
        kinetyka nie zmienia jakościowo kształtu.
\end{itemize}
To sugeruje, że $\mathcal{T}_b$ zbiega do unikalnego punktu stałego
w sektorze kinetycznym, co jest konieczne dla OP-2.

\paragraph{Argument Q-II: Twierdzenie de Giorgi–Nash–Moser (analityczny).}
Niech $\Phi_B(x)$ będzie polem blokowym \eqref{eq:block_avg}.
Jeżeli działanie efektywne $S_\mathrm{eff}[\Phi_B]$ ma regularny
człon gradientowy (co gwarantuje $K(0) = 0$ i $K(\Phi) > 0$ dla $\Phi > 0$),
to z twierdzenia de Giorgi–Nash–Moser \cite{DeGiorgi1957}:
\begin{equation}
  \|\Phi_B\|_{C^{0,\alpha}(\Omega)} \leq C \|\Phi_B\|_{L^2(\Omega)},
  \quad \alpha > 0,
\end{equation}
tj. $\Phi_B$ jest funkcją Höldera ciągłą, a więc zbiega do funkcji ciągłej
$\Phi \in C^{0,\alpha}$ przy $b \to \infty$.

\paragraph{Argument Q-III: Teoria homogenizacji (analityczny, szkic).}
Dla periodycznych (lub quasi-periodycznych) substratów istnieje teoria
homogenizacji \cite{Bensoussan1978}, która gwarantuje, że
$\Phi_B(x) \rightharpoonup \Phi(x)$ słabo w $H^1(\Omega)$,
gdzie $\Phi$ spełnia równanie jednorodne:
\begin{equation}
  -\nabla \cdot (K_\mathrm{hom}(\Phi)\nabla\Phi) + V'_\mathrm{hom}(\Phi) = q\rho.
\end{equation}
Identyfikacja $K_\mathrm{hom} = K_\mathrm{TGP}$ i $V'_\mathrm{hom} = U'_\mathrm{TGP}$
jest \emph{warunkiem zamknięcia OP-2}.

\subsection*{Q.4 Skala korelacji i relacja $a_\Gamma \cdot \Phi_0 = 1$}

\begin{propozycja}[Φ₀ z przejścia fazowego substratu]
\label{prop:phi0_from_Tc}
W fazie $T < T_c$ (spontane łamanie $\mathbb{Z}_2$):
\begin{equation}
  \langle\phi^2\rangle = v^2(T) > 0,
  \qquad v^2(T_c) = 0.
\end{equation}
Tło próżniowe TGP jest naturalnie:
\begin{equation}
  \boxed{\Phi_0 = \langle\phi^2\rangle\big|_{T \ll T_c} = v^2(T_\mathrm{eff})},
  \label{eq:phi0_from_TC}
\end{equation}
gdzie $T_\mathrm{eff}$ jest efektywną temperaturą substratu przy
obecnym wieku Wszechświata.
\end{propozycja}

\begin{propozycja}[Hipoteza $a_\Gamma \cdot \Phi_0 = 1$ jako warunek samospójności]
\label{prop:agamma_phi0}
Definiując skalę korelacji substratu:
\begin{equation}
  \xi_\mathrm{corr}(T) = a_\mathrm{sub} \cdot |1 - T/T_c|^{-\nu},
  \quad \nu \approx 0{,}63 \text{ (3D Ising)},
\end{equation}
i identyfikując $a_\Gamma \equiv 1/\xi_\mathrm{corr}$, hipoteza
$a_\Gamma \cdot \Phi_0 = 1$ jest równoznaczna z:
\begin{equation}
  \xi_\mathrm{corr} = \Phi_0
  \quad\Longleftrightarrow\quad
  a_\mathrm{sub} \cdot |1 - T_\mathrm{eff}/T_c|^{-\nu} = v^2(T_\mathrm{eff}).
  \label{eq:self_consistency}
\end{equation}
Warunek \eqref{eq:self_consistency} jest \emph{równaniem samospójnym}
wyznaczającym $T_\mathrm{eff}$ z parametrów substratu $(J, a_\mathrm{sub}, T_c)$.
\end{propozycja}

\begin{uwaga}[Weryfikacja numeryczna]
Skrypt \texttt{tgp\_agamma\_phi0\_test.py} sprawdza, czy $a_\Gamma \cdot \Phi_0 = 1$
jest spełnione numerycznie z trzech niezależnych ścieżek obserwacyjnych:
\begin{enumerate}
  \item $\Phi_0(\Lambda) = 96\pi G_0 \rho_\Lambda / H_0^2 \approx 23{,}3$
  \item $\Phi_0(r_{21}) = (\alpha_K \sqrt{r_{21}})^{3/5} \approx 25{,}4$
  \item $\Phi_0(\kappa) = 3/(4\kappa_\mathrm{obs}) \approx 24{,}7$
\end{enumerate}
Rozrzut $< 10\%$ (patrz ANALIZA\_SPOJNOSCI\_v3.md), a DESI DR2 mierzy
$a_\Gamma \cdot \Phi_0 = 1{,}005 \pm 0{,}005$.
\end{uwaga}

\subsection*{Q.5 Mapa otwartego problemu OP-2}

\begin{center}
\begin{tabular}{llll}
\toprule
Krok & Opis & Typ & Status \\
\midrule
CG-1 & Istnienie i jedyność $S^*$ (Banach) & [AN] & \textbf{[OTWARTY]} \\
CG-2 & LPA' z $K_k(\rho)$, zbieżność do $K_\mathrm{TGP}$ & [NUM] & \textbf{[ZAMKNIĘTY NUM]} \\
CG-3 & Zbieżność $\Phi_B \to \Phi$ w $H^1$ (homogenizacja) & [AN] & \textbf{[OTWARTY]} \\
CG-4 & Identyfikacja $K_\mathrm{hom} = K_\mathrm{TGP}$ & [AN+NUM] & \textbf{[OTWARTY]} \\
CG-5 & Φ₀ z równania samospójnego \eqref{eq:self_consistency} & [NUM] & \textbf{[ZAMKNIĘTY NUM]} \\
\bottomrule
\end{tabular}
\end{center}

\medskip\noindent
Krok CG-2 (LPA') jest \textbf{zamknięty numerycznie} skryptem
\texttt{tgp\_erg\_lpa\_prime.py} (8/8 PASS, 2026-04-02):
$\rho_0^*=0.03045$, $\nu=0.749$, $K_\mathrm{IR}/K_\mathrm{UV}=1.000$,
$v_2=0.06090$.
Krok OP-1 (LPA' z $\eta\neq0$) zamknięty skryptem
\texttt{tgp\_erg\_eta\_lpa\_prime.py} (10/10 PASS):
$\eta^*=0.04419$ (samospójnie).
Kroki CG-1, CG-3, CG-4 wymagają zaawansowanej matematyki
(ERG funkcjonalne, teoria homogenizacji de Giorgi, analiza funkcjonalna).
Krok CG-5 zamknięty skryptem \texttt{tgp\_cg5\_phi0\_self\_consistency.py}
(8/8 PASS, 2026-04-02): $T^*/T_c = 0.7955$, $a_\mathrm{sub,fiz} = 9.072$,
$a_\Gamma\cdot\Phi_0 = 1.000000$.

\subsection*{Q.6 Program numeryczny dla CG-2 i CG-5}

Skrypt \texttt{tgp\_erg\_lpa\_prime.py} (zastąpił
\texttt{tgp\_erg\_wetterich.py} w roli głównego skryptu CG-2) realizuje:
\begin{enumerate}
  \item Pełne równanie LPA' dla $K_k(\rho)$ (eq.~N.6 w Dodatku~N):
        \begin{equation}
          \partial_t K_k(\rho) = \frac{k^5}{6\pi^2} \cdot
          \frac{(K_k' + 2\rho K_k'') - K_k D_k}{(k^2 + U_k' + 2\rho U_k'')^2},
        \end{equation}
        gdzie $D_k$ zawiera wyższe pochodne $U_k$.
  \item Test zbieżności: czy $K_k(\rho)/k^\eta \to K_\mathrm{IR}(\rho)$ przy $k \to 0$?
  \item Wyznaczenie $v^2 = \langle\phi^2\rangle$ z profilu IR: $v^2 = 2\rho_0^*$
        gdzie $\rho_0^*$ jest minimum punktu stałego $\tilde u^*(\tilde\rho)$.
  \item Weryfikacja samospójności: $a_\Gamma \cdot 2\rho_0^* \stackrel{?}{=} 1$.
\end{enumerate}

\subsection*{Q.7 Implikacje dla TGP}

Jeżeli OP-2 zostanie domknięty (choćby numerycznie przez CG-2):
\begin{itemize}
  \item $\Phi_0$ staje się \textbf{predykcją} z parametrów substratu $(J, T_c)$
        — nie wolnym parametrem. Redukuje $N_\mathrm{param}$ z 2 do 1.
  \item Relacja $a_\Gamma \cdot \Phi_0 = 1$ zyskuje \textbf{mechanizm}
        (warunek samospójności, nie przypadkowa zbieżność).
  \item Iloraz $M_\mathrm{obs}/N_\mathrm{param} \geq 12$ — wszystkie
        obserwable predykcyjne przy jednym wolnym parametrze.
\end{itemize}

\bigskip
\begin{center}
\fbox{\parbox{0.85\textwidth}{
\textbf{Podsumowanie OP-2.}
Problem gruboziarnienia $\Gamma \to \Phi$ jest precyzyjnie sformułowany
jako zbieżność operatora blokowania $\mathcal{T}_b$ do punktu stałego $S^*$.
Trzy argumenty (ERG, de Giorgi, homogenizacja) wspierają zbieżność,
ale żaden z nich nie stanowi pełnego dowodu w kontekście TGP.
Zamknięcie CG-2 (numeryczne LPA') jest celem najbliższym dostępnemu
numerycznie. Pe\l{}ny dow\'od analityczny to zadanie na 6--12 miesi\k{e}cy.
}}
\end{center}

% ==============================================================
% Q.7 Propozycje domknięcia continuum (v47b)
% ==============================================================
\subsection*{Q.7 Propozycje domkni\k{e}cia continuum (v47b)}
\label{ssec:Q-continuum-proposals}

Poni\.zsze dwie propozycje s\k{a} \textbf{kandydatami na formalne mosty},
nie pe\l{}nymi twierdzeniami. Oznaczone jako \statuslabel{Szkic}.

\begin{proposition}[S\l{}abe twierdzenie continuum]%
\label{prop:Q-weak-continuum}
\textbf{Status: \statuslabel{Szkic}.}
Niech $\Phi_B(x) = N_B^{-1}\sum_{i \in B(x)} \phi_i^2$
b\k{e}dzie polem blokowym na substracie $\Gamma$
z~lokalnym hamiltonianem i~d\l{}ugo\'sci\k{a} korelacji $\xi_{\rm corr}$.
Je\.zeli:
\begin{enumerate}[nosep]
  \item $\sup_B \|\Phi_B\|_{L^2(K)} < \infty$ dla ka\.zdego
    zwartego~$K$,
  \item sektor gradientowy energii swobodnej spe\l{}nia
    $F_B[\Phi_B] \ge c_* \|\nabla\Phi_B\|^2_{L^2(K)} - C_K$
    z~$c_* > 0$,
  \item korelacje dwupunktowe malej\k{a} szybciej ni\.z
    $\exp(-r/\xi_{\rm corr})$,
\end{enumerate}
to rodzina $\{\Phi_B\}$ jest prezwarta s\l{}abo w~$H^1_{\rm loc}$
i~mocno w~$L^2_{\rm loc}$ (twierdzenie Rellicha--Kondrachova).
Punkty skupienia spe\l{}niaj\k{a} r\'ownanie Eulera--Lagrange'a
pewnego lokalnego funkcjona\l{}u efektywnego.

\medskip\noindent
\textit{Znaczenie}: zast\k{e}puje niedost\k{e}pny pe\l{}ny dow\'od
(CG-1 Banacha) s\l{}abszym, ale publikowalnym rezultatem
o~istnieniu granicy ci\k{a}g\l{}ej podsekwencji.
Najtrudniejszy punkt techniczny to uzasadnienie
dodatnio\'sci efektywnej sztywno\'sci po blokowaniu.
\end{proposition}

\begin{proposition}[$\alpha_{\rm eff} = 2$ z~przej\'scia
  $\phi \to \Phi = \phi^2$]%
\label{prop:Q-alpha-from-phi-squared}
\textbf{Status: \statuslabel{Propozycja}.}
Je\.zeli $\Phi = \phi^2$ z~$\phi > 0$, to:
\begin{align}
  \nabla\Phi &= 2\phi\,\nabla\phi, &
  (\nabla\Phi)^2 &= 4\Phi\,(\nabla\phi)^2, &
  (\nabla\phi)^2 &= \frac{(\nabla\Phi)^2}{4\Phi}.
\end{align}
Ka\.zdy regularny cz\l{}on gradientowy dla amplitudy,
\[
  \int Z(\phi)\,(\nabla\phi)^2\,d^3x,
\]
po zapisaniu w~zmiennej $\Phi$ przyjmuje posta\'c
\[
  \int \frac{Z(\sqrt{\Phi})}{4\Phi}\,(\nabla\Phi)^2\,d^3x.
\]
Je\.zeli $Z(\phi) \sim \phi^2$ (co sugeruje struktura substratowa),
to $Z(\sqrt{\Phi})/\Phi = \mathrm{const} + O(\Phi - \Phi_0)$,
a~wariacja generuje dok\l{}adnie cz\l{}on $(\nabla\Phi)^2/\Phi$,
czyli \textbf{$\alpha_{\rm eff} = 2$}.

\medskip\noindent
\textit{Wniosek}: $\alpha = 2$ w~r\'ownaniu pola TGP
\textbf{nie jest dowolnym wsp\'o\l{}czynnikiem}, lecz
algebraiczn\k{a} konsekwencj\k{a} przej\'scia od amplitudy
do jej kwadratu. To mocno wspiera naturalnie\'s\'c r\'ownania pola.
\end{proposition}
